% ch59_toward_hott.tex — Volume IV Chapter 23
% Retrofitted with full Vol I-III pedagogy: cycle stages, etymology, dialogs.

\chapter{Toward Homotopy Type Theory}

\begin{overviewbox}
\term{Homotopy Type Theory} (HoTT) is the intersection of $\infty$-category
theory and dependent type theory. Its central insight is the
\term{univalence axiom}: equivalent types are equal — formally,
$(A \simeq B) \simeq (A = B)$ as types. Univalence is the type-theoretic
avatar of the Yoneda philosophy, and the resulting framework supplies a
foundational language for mathematics in which:
\begin{itemize}
  \item Spaces (more accurately: $\infty$-groupoids) are primitive
        rather than derived from sets.
  \item Identity is structural: two structures with all the same
        properties are formally equal.
  \item Constructive proofs and computer-checkable verification become
        natural.
\end{itemize}

This chapter is a preview, parallel to Chapter~58 on DAG. We introduce
the basics — dependent types, identity types, the univalence axiom,
$\infty$-topos models — and gesture at where HoTT is heading: the
\term{univalent foundations} program (Voevodsky), modal type theories
(synthetic homotopy theory), and applications in algebraic topology and
constructive mathematics.

Full development would require a dedicated volume; the canonical
references are \emph{The HoTT Book} (Univalent Foundations Project,
2013) and Riehl–Shulman, \emph{Synthetic theory of $\infty$-categories
inside Homotopy Type Theory}. After this chapter, the reader can engage
those works.
\end{overviewbox}


% ═══════════════════════════════════════════════════════════════════════════
\section{Dependent Type Theory: A Refresher}
\label{sec:dtt-refresher}
% ═══════════════════════════════════════════════════════════════════════════

\subsection{Informally}

A type theory has \emph{types} $A, B, \ldots$ and \emph{terms} (also
called \emph{elements}) of those types, written $a : A$ to mean ``$a$
is a term of type $A$.'' Function types $A \to B$ have terms that are
functions; the \term{lambda calculus} provides the syntax.

\term{Dependent type theory} allows types to depend on terms: $B(x)$
might be a type whose definition depends on $x : A$. Two key constructions:
\begin{itemize}
  \item \term{Dependent function type} $\prod_{x : A} B(x)$: functions
        taking $x : A$ and producing a term of $B(x)$.
  \item \term{Dependent pair type} $\sum_{x : A} B(x)$: pairs $(x, y)$
        with $x : A$ and $y : B(x)$.
\end{itemize}
These generalize $A \to B$ (when $B$ is constant) and $A \times B$
(when $B$ is constant).

\subsection{Formalizing}

\begin{nameorigin}
\term{Type theory} dates to \emph{Bertrand Russell}'s 1903 attempt to
resolve set-theoretic paradoxes by stratifying mathematical objects
into a hierarchy of types. The modern computational version is
\emph{Alonzo Church}'s simply-typed lambda calculus (1940); the
\emph{dependent} extension is due to \emph{Per Martin-Löf} (1972),
adding types that depend on terms. The Curry-Howard-Lambek correspondence
(see Vol~I Ch.~46 / Vol~V) identifies type-theoretic proofs with
categorical morphisms, with dependent types corresponding to indexed
products/coproducts. The $\infty$-categorical Martin-Löf type theory
(Univalent Foundations, HoTT) is the modern synthesis.
\end{nameorigin}

\begin{definition}[Dependent type theory, informally]
\label{def:dtt}
A \term{dependent type theory} consists of:
\begin{itemize}
  \item A judgment of typing: $x : A \vdash B \, \mathrm{type}$.
  \item Function types $\prod_{x : A} B(x)$ with lambda-abstraction
        and application.
  \item Sigma types $\sum_{x : A} B(x)$ with pairing and projections.
  \item Universe types $\mathcal{U}, \mathcal{U}_1, \mathcal{U}_2, \ldots$
        with $\mathcal{U}_i : \mathcal{U}_{i+1}$ (avoiding Russell's
        paradox).
  \item Various other type formers: inductive types, $\mathbb{N}, \mathbb{B}$,
        etc.
\end{itemize}
\end{definition}

\begin{remark}[Curry–Howard–Lambek]
The classical \term{Curry–Howard–Lambek correspondence} (Chapter~46 / FP
volume) identifies:
\begin{itemize}
  \item Types $\leftrightarrow$ propositions $\leftrightarrow$ objects
        in a CCC.
  \item Terms $\leftrightarrow$ proofs $\leftrightarrow$ morphisms.
  \item $A \to B$ $\leftrightarrow$ implication $\leftrightarrow$ exponential.
  \item $A \times B$ $\leftrightarrow$ conjunction $\leftrightarrow$ product.
  \item $A + B$ $\leftrightarrow$ disjunction $\leftrightarrow$ coproduct.
\end{itemize}
Dependent type theory extends this to higher-categorical settings:
$\prod, \sum$ correspond to indexed product/coproduct $=$ dependent
adjoints.
\end{remark}


% ═══════════════════════════════════════════════════════════════════════════
\section{Identity Types and Path Types}
\label{sec:identity-types}
% ═══════════════════════════════════════════════════════════════════════════

\subsection{Formalizing}

\begin{nameorigin}
\term{Identity type} was introduced by \emph{Per Martin-Löf} (1972) as
part of his intuitionistic type theory. The classical view: $a =_A b$
is a type whose inhabitants are \emph{proofs of equality} between $a$
and $b$. The radical move of HoTT (Awodey-Warren 2009, Voevodsky $\sim$
2010): take seriously the structure of identity types as
\emph{$\infty$-groupoids}. The path interpretation — $a = b$ is the
mapping space $\mathrm{Path}(a, b)$ — connects type theory to homotopy
theory in a way that has reshaped foundations.
\end{nameorigin}

\begin{definition}[Identity type]
\label{def:identity-type}
For terms $a, b : A$, the \term{identity type} $a =_A b$ (or simply
$a = b$) is a type whose terms witness $a$ and $b$ being identical.
\begin{itemize}
  \item Introduction: $\mathrm{refl}_a : a = a$ for every $a : A$.
  \item Induction (J-rule): properties holding on $\mathrm{refl}$ extend
        to all of $a = b$.
\end{itemize}
\end{definition}

\begin{remark}[Path interpretation]
In the homotopy-theoretic interpretation, $A$ is an $\infty$-groupoid
(Kan complex), terms $a, b : A$ are vertices, and $a = b$ is the
\emph{mapping space} $\mathrm{Path}_A(a, b)$. $\mathrm{refl}_a$ is the
constant path at $a$.

Iterating: $\mathrm{refl}_a = \mathrm{refl}_a$ is the type of $2$-paths
(homotopies between paths); and so on. The infinite tower of identity
types is exactly the infinite tower of higher cells in an
$\infty$-groupoid.
\end{remark}


% ═══════════════════════════════════════════════════════════════════════════
\section{The Univalence Axiom}
\label{sec:univalence}
% ═══════════════════════════════════════════════════════════════════════════

\subsection{Formalizing}

\begin{definition}[Equivalence of types]
\label{def:equiv-types}
For types $A, B$, the type $A \simeq B$ of \term{equivalences} from $A$
to $B$ consists of functions $f : A \to B$ together with proofs that $f$
has both a left and a right homotopy inverse.

(The technical detail: equivalences are functions whose fibres are
contractible, equivalent to ``invertible up to coherent homotopy.'')
\end{definition}

\begin{nameorigin}
The \term{univalence axiom} was introduced by \emph{Vladimir Voevodsky}
(2010) as part of his \emph{Univalent Foundations} program. Voevodsky
(1966--2017) was a Fields Medalist (2002) primarily known for motivic
cohomology, and in his final decade turned to foundations of mathematics
in HoTT. The univalence axiom — that equivalent types are equal — is
the formal expression of the structural identity principle: ``isomorphic
structures should be treated as identical.'' The HoTT Book (Univalent
Foundations Project 2013) is the canonical reference.
\end{nameorigin}

\begin{theorem}[Univalence axiom, Voevodsky]
\label{thm:univalence}
For all types $A, B : \mathcal{U}$, the canonical map
\begin{equation*}
  \mathrm{idtoeq} : (A = B) \to (A \simeq B)
\end{equation*}
is itself an equivalence. In particular, \term{equivalent types are
equal}.
\end{theorem}

\begin{remark}[Univalence as Yoneda]
\label{rem:univalence-yoneda}
Univalence is the type-theoretic analogue of the $\infty$-categorical
Yoneda philosophy: ``an object is determined by what maps in.'' In HoTT:
``a type is determined by what is equivalent to it.''

At the model level: HoTT is interpreted in $\infty$-toposes, where the
universe $\mathcal{U}$ classifies $\infty$-groupoids; univalence
corresponds to the (small-object) representability of the universe
fibration.
\end{remark}

\begin{example_thm}[Function extensionality from univalence]
\label{ex:funext}
A consequence of univalence: \term{function extensionality}. If $f, g :
A \to B$ have $(x : A) \vdash f(x) = g(x)$, then $f = g$.

Function extensionality is not provable in vanilla dependent type theory
but follows from univalence. This is one of the deep payoffs.
\end{example_thm}


% ═══════════════════════════════════════════════════════════════════════════
\section{$\infty$-Topos Models}
\label{sec:infty-topos}
% ═══════════════════════════════════════════════════════════════════════════

\subsection{Formalizing}

\begin{nameorigin}
\term{$\infty$-topos} extends \emph{Grothendieck topos} (1957--1963) to
the homotopy-coherent setting. A classical Grothendieck topos is a
category of sheaves on a site; an $\infty$-topos is a category of
$\infty$-sheaves on an $\infty$-site, with the sheaf condition phrased
in $\infty$-categorical descent. The development is largely due to
\emph{Lurie} (HTT Ch.~6), building on \emph{Charles Rezk}, \emph{Bertrand
Toën}, and others. The semantic claim: $\infty$-toposes are exactly
the categories of mathematical structures definable in HoTT.
\end{nameorigin}

\begin{definition}[$\infty$-topos]
\label{def:infty-topos}
An \term{$\infty$-topos} is an accessible left-exact localization of a
presheaf $\infty$-category $\mathcal{P}(\mathcal{C})$. Equivalently:
an $\infty$-category $\mathcal{X}$ satisfying Giraud's axioms in the
$\infty$-categorical setting (colimit-presentable plus universality
properties).
\end{definition}

\begin{theorem}[Infty-topos models of HoTT]
\label{thm:infty-topos-hott}
Every $\infty$-topos $\mathcal{X}$ provides a model for HoTT: types
are objects of $\mathcal{X}$; dependent types are morphisms; identity
types are diagonals; univalence holds because $\mathcal{X}$'s universe
classifies $\infty$-groupoid-valued sheaves.

The principal example is $\mathcal{S}$ itself, which models HoTT
trivially (types $=$ $\infty$-groupoids).
\end{theorem}

\begin{remark}[Lurie's HTT and HoTT]
The model-theoretic foundations of $\infty$-toposes (Lurie, HTT §6) and
the type-theoretic foundations of HoTT (Univalent Foundations) are
remarkably compatible: $\infty$-toposes give the semantics that HoTT
specifies syntactically. The synthetic and analytic pictures meet here.
\end{remark}


% ═══════════════════════════════════════════════════════════════════════════
\section{Reading Recommendations}
\label{sec:hott-reading}
% ═══════════════════════════════════════════════════════════════════════════

\begin{remark}[Where to go]
Foundational references for HoTT:
\begin{itemize}
  \item \emph{The HoTT Book} (\textbf{Univalent Foundations Project}, 2013):
        the canonical introduction. Available at \texttt{homotopytypetheory.org}.
  \item Riehl–Shulman, \emph{A type theory for synthetic $\infty$-categories}:
        $\infty$-category theory \emph{within} HoTT — a more advanced
        treatment.
  \item Voevodsky's \emph{Univalent Foundations Project} writings.
  \item Awodey–Warren, \emph{Homotopy theoretic models of identity types}:
        the original semantic insight.
\end{itemize}

For HoTT applied to algebraic topology:
\begin{itemize}
  \item Brunerie's PhD thesis on $\pi_4(S^3)$.
  \item Buchholtz–van Doorn–Rijke, \emph{Higher groups in homotopy type theory}.
  \item Various papers on synthetic homotopy theory.
\end{itemize}
\end{remark}


% ═══════════════════════════════════════════════════════════════════════════
\section{Exercises}
\label{sec:hott-exercises}
% ═══════════════════════════════════════════════════════════════════════════

\begin{exercisesbox}
\begin{qa}
\qaitem{What is dependent type theory?}{A type theory where types can depend on terms: $x : A \vdash B(x) \, \mathrm{type}$. Includes dependent function types $\prod_{x : A} B(x)$ and dependent pair types $\sum_{x : A} B(x)$.}
\qaitem{What is the Curry-Howard-Lambek correspondence?}{Types $\leftrightarrow$ propositions $\leftrightarrow$ objects in a CCC. Terms $\leftrightarrow$ proofs $\leftrightarrow$ morphisms. Function types $\leftrightarrow$ exponentials. Connects logic, type theory, and category theory.}
\qaitem{Define the identity type.}{For $a, b : A$, $a =_A b$ is a type with $\mathrm{refl}_a : a = a$ as canonical inhabitant. Induction principle: properties on $\mathrm{refl}$ extend.}
\qaitem{What's the path interpretation?}{$A$ is an $\infty$-groupoid; $a = b$ is the mapping space $\mathrm{Path}_A(a, b)$. Iterating gives the full tower of higher cells.}
\qaitem{Define equivalence of types.}{$A \simeq B$ is the type of equivalences: functions with contractible homotopy fibres. Equivalent to ``invertible up to coherent homotopy.''}
\qaitem{State the univalence axiom.}{$\mathrm{idtoeq} : (A = B) \to (A \simeq B)$ is itself an equivalence. Equivalent types are equal.}
\qaitem{What does univalence ``do''?}{It identifies type-level equivalence with type-level equality, so reasoning about equivalent structures is the same as reasoning about equal ones. Reverses the classical ``identity'' problem of foundations.}
\qaitem{What's function extensionality?}{$(\forall x, f(x) = g(x)) \Rightarrow f = g$. Provable from univalence; independent of vanilla dependent type theory.}
\qaitem{How does univalence relate to Yoneda?}{Both say: objects determined by relationships. Yoneda: object $=$ representable functor. Univalence: type $=$ equivalence class under $\simeq$. Same philosophical content, different formal setup.}
\qaitem{Define an $\infty$-topos.}{An accessible left-exact localization of $\mathcal{P}(\mathcal{C})$ for small $\mathcal{C}$. Equivalent: $\infty$-cat satisfying $\infty$-Giraud's axioms.}
\qaitem{What's the principal $\infty$-topos model of HoTT?}{$\mathcal{S}$ itself: types $=$ $\infty$-groupoids, types-of-types $=$ universe in $\mathcal{S}$, identity types $=$ path spaces.}
\qaitem{Why is $\mathcal{S}$ an $\infty$-topos?}{$\mathcal{S} = \mathcal{P}(\mathrm{pt}) = \mathrm{Fun}(\mathrm{pt}^{\mathrm{op}}, \mathcal{S}) = \mathcal{S}$ trivially; it satisfies the $\infty$-Giraud axioms.}
\qaitem{What's an $\infty$-topos's universe?}{A subobject classifier in the $\infty$-categorical sense: an object $\mathcal{U}$ such that morphisms $X \to \mathcal{U}$ classify (small) subobjects of $X$. Generalizes Set's universe to $\infty$-groupoids.}
\qaitem{Why does univalence hold in $\infty$-topos models?}{Because the universe's classifying property includes equivalences: a small subobject up to equivalence is the same as a small subobject up to equality.}
\qaitem{What is synthetic homotopy theory?}{Doing homotopy theory inside HoTT, without recourse to topological spaces or simplicial sets. Computations of $\pi_n(S^k)$ can be done synthetically.}
\qaitem{Has $\pi_4(S^3)$ been computed in HoTT?}{Yes — Guillaume Brunerie's PhD thesis (2016) gives a constructive proof that $\pi_4(S^3) = \mathbb{Z}/2$. Verified by computer.}
\qaitem{What's higher inductive types (HITs)?}{Type formers that introduce both points and higher-dimensional paths. Examples: $\mathbb{S}^1$ as a HIT with one point and one path-loop; sphere $S^n$ via iterated HITs.}
\qaitem{Why is HoTT a candidate foundation?}{Because it captures higher-categorical phenomena natively, supports computer verification, and makes structural identity ($\simeq$ $=$ $=$) a primitive. Voevodsky's Univalent Foundations project.}
\qaitem{Is HoTT consistent?}{Yes, in the sense that it has set-theoretic models ($\infty$-toposes). Like ZFC, consistency requires a meta-theory.}
\qaitem{Where does HoTT meet $\infty$-category theory?}{They're two sides of the same coin: HoTT is the internal language of $\infty$-toposes; $\infty$-category theory is the external/semantic theory of HoTT.}
\qaitem{What about modal HoTT?}{Adds modal operators (necessity, possibility) to type theory, used to model synthetic differential geometry, cohesion, and physical theories.}
\qaitem{What's the next (final) chapter?}{Ch~60: The Higher Landscape — synthesizing Vol~IV, mapping its chapters to standard literature, and previewing a potential Vol~V on applications.}
\qaitem{What's Vol~V (if written) supposed to be?}{The fulfillment of the original v0.5 promise: connections to other branches of mathematics. Algebra (derived/triangulated), topology (spectra), logic (toposes), CS (type theory). All deferred to ``after Vol~IV.''}
\end{qa}
\end{exercisesbox}


% ═══════════════════════════════════════════════════════════════════════════
\section*{Chapter Summary}
\addcontentsline{toc}{section}{Chapter Summary}
% ═══════════════════════════════════════════════════════════════════════════

\begin{summarybox}
\textbf{Dependent type theory.}\quad Types depending on terms;
$\prod_{x:A} B(x)$ and $\sum_{x:A} B(x)$ as dependent function and pair
types. Curry-Howard-Lambek extended to higher categories.

\textbf{Identity types.}\quad $a =_A b$ as the type of identifications.
Path interpretation: $a = b$ is the mapping space in the $\infty$-groupoid $A$.

\textbf{Univalence.}\quad $(A = B) \simeq (A \simeq B)$. Equivalent types
are equal. The structural identity principle made formal. Function
extensionality follows.

\textbf{$\infty$-topos models.}\quad Every $\infty$-topos models HoTT;
$\mathcal{S}$ is the principal example. Synthetic and semantic
foundations agree.

\textbf{Reading.}\quad The HoTT Book; Riehl-Shulman; Voevodsky's
Univalent Foundations; Brunerie on $\pi_4(S^3)$.

\textbf{Vol IV synthesis closes Chapter~60.}\quad The original v0.5
promise (connections to other math fields) becomes potential Vol~V.
\end{summarybox}


% ═══════════════════════════════════════════════════════════════════════════
\section*{Formal Restatement}
\addcontentsline{toc}{section}{Formal Restatement}
% ═══════════════════════════════════════════════════════════════════════════

\begin{formalrestatement}
\textbf{Dependent type theory.}\quad Types $A : \mathcal{U}$;
dependent types $x : A \vdash B(x) : \mathcal{U}$;
$\prod_{x : A} B(x)$ and $\sum_{x : A} B(x)$.

\medskip
\textbf{Identity type.}\quad $a, b : A \Rightarrow (a = b) : \mathcal{U}$;
$\mathrm{refl}_a : a = a$; J-rule eliminator.

\medskip
\textbf{Equivalence.}\quad $A \simeq B = \sum_{f : A \to B} \mathrm{isEquiv}(f)$.

\medskip
\textbf{Univalence axiom.}\quad
$\mathrm{idtoeq}_{A, B} : (A = B) \to (A \simeq B)$ is an equivalence.

\medskip
\textbf{$\infty$-topos.}\quad Accessible left-exact localization of
$\mathcal{P}(\mathcal{C})$; equivalent: $\infty$-Giraud's axioms.

\medskip
\textbf{Model theorem.}\quad Every $\infty$-topos models univalent
dependent type theory. $\mathcal{S}$ is the principal example.
\end{formalrestatement}
